% LightGBM reduces split-search work along threshold, sample, and feature axes.
% Schematic only; marks and sparsity patterns encode no measured values.
% Canvas: 164 mm x 64 mm. Dependencies: project palette and TikZ.
\begin{tikzpicture}[
  x=1mm,y=1mm,
  font=\footnotesize,
  text=BookInk,
  flow/.style={
    -{Latex[length=2.2mm,width=1.5mm]},
    draw=BookMuted,
    line width=0.8pt
  },
  keep/.style={
    circle,
    draw=FigureLoss,
    fill=FigureLoss!12,
    line width=0.8pt,
    minimum size=5mm,
    inner sep=0pt
  },
  sample/.style={
    circle,
    draw=FigureInput,
    fill=FigureInput!8,
    line width=0.7pt,
    minimum size=4mm,
    inner sep=0pt
  }
]
  \path[use as bounding box] (0,0) rectangle (164,64);

  \node[font=\heiti\small] at (27,59) {(a) 直方图：减少候选阈值};
  \draw[BookMuted,line width=0.7pt] (5,35)--(49,35);
  \foreach \x in {8,12,15,18,22,24,28,31,35,39,42,46} {
    \fill[FigureInput] (\x,35) circle[radius=1.1mm];
  }
  \foreach \x in {5,16,27,38,49} {
    \draw[FigureModel,line width=0.8pt] (\x,28)--(\x,43);
  }
  \node[text=BookMuted] at (27,24) {连续取值映射到有限个桶};
  \node[draw=FigureOutput,fill=FigureOutput!6,line width=0.8pt,
    minimum width=38mm,minimum height=9mm,align=center] (bins) at (27,10)
    {只扫描桶边界\\并累计梯度统计};
  \draw[flow] (27,22)--(bins);

  \draw[BookRule,line width=0.6pt] (54.5,3)--(54.5,61);

  \node[font=\heiti\small] at (82,59) {(b) GOSS：减少参与统计的样本};
  \node[anchor=east,text=BookMuted] at (72,47) {较大的$|g_i|$};
  \foreach \x in {78,86,94,102} {
    \node[keep] at (\x,47) {};
  }
  \node[anchor=east,text=BookMuted] at (72,34) {较小的$|g_i|$};
  \foreach \x in {75,81,87,93,99,105} {
    \node[sample] at (\x,34) {};
  }
  \node[draw=FigureOutput,fill=FigureOutput!6,line width=0.8pt,
    minimum width=42mm,minimum height=15mm,text width=38mm,align=center]
    (goss) at (82,12)
    {保留$|g_i|$较大的样本\\随机抽取并重加权部分$|g_i|$较小的样本};
  \draw[flow,draw=FigureLoss] (105,47)--(108,47)|-(goss.east);
  \draw[flow] (87,31.5)--(87,27)--(57,27)|-(goss.west);
  \draw[BookMuted,dashed,line width=0.7pt] (81,31)--(81,37);
  \draw[BookMuted,dashed,line width=0.7pt] (93,31)--(93,37);
  \draw[BookMuted,dashed,line width=0.7pt] (105,31)--(105,37);

  \draw[BookRule,line width=0.6pt] (109.5,3)--(109.5,61);

  \node[font=\heiti\small] at (137,59) {(c) EFB：减少有效特征};
  \node[anchor=east,text=BookMuted] at (121,45) {$x_a$};
  \node[anchor=east,text=BookMuted] at (121,35) {$x_b$};
  \foreach \x in {125,131,137,143,149,155} {
    \draw[BookRule,line width=0.6pt] (\x-2.2,42.5) rectangle (\x+2.2,47.5);
    \draw[BookRule,line width=0.6pt] (\x-2.2,32.5) rectangle (\x+2.2,37.5);
  }
  \fill[FigureInput] (122.8,42.5) rectangle (127.2,47.5);
  \fill[FigureInput] (140.8,42.5) rectangle (145.2,47.5);
  \fill[FigureInput] (128.8,32.5) rectangle (133.2,37.5);
  \fill[FigureInput] (146.8,32.5) rectangle (151.2,37.5);
  \node[text=BookMuted,align=center] at (139,25)
    {非零位置很少冲突};
  \node[draw=FigureOutput,fill=FigureOutput!6,line width=0.8pt,
    minimum width=42mm,minimum height=10mm,text width=38mm,align=center]
    (bundle) at (137,10)
    {编码为一个捆绑特征};
  \draw[flow] (139,21)--(bundle);
\end{tikzpicture}
